\documentclass[11pt,a4paper]{article}

% ─── Packages ───────────────────────────────────────────────────────
\usepackage[utf8]{inputenc}
\usepackage[T1]{fontenc}
\usepackage{lmodern}
\usepackage[margin=1in]{geometry}
\usepackage{amsmath,amssymb,amsthm}
\usepackage{booktabs}
\usepackage{graphicx}
\usepackage{listings}
\usepackage{xcolor}
\usepackage{hyperref}
\usepackage{cleveref}
\usepackage{algorithm}
\usepackage{algpseudocode}
\usepackage{enumitem}
\usepackage{float}
\usepackage{caption}
\usepackage{subcaption}
\usepackage{tikz}
\usetikzlibrary{arrows.meta,positioning,shapes.geometric,fit,calc}
\usepackage{textcomp}

% ─── Compatibility: define semantic brackets without stmaryrd ──────
\providecommand{\llbracket}{[\![}
\providecommand{\rrbracket}{]\!]}

% ─── Listings style for ChimeraLang ────────────────────────────────
\definecolor{chimkey}{RGB}{0,100,180}
\definecolor{chimtype}{RGB}{160,32,80}
\definecolor{chimstr}{RGB}{0,128,64}
\definecolor{chimcom}{RGB}{128,128,128}
\definecolor{chimbg}{RGB}{248,248,252}

\lstdefinelanguage{chimera}{
  keywords={val,fn,gate,goal,reason,about,end,must,allow,forbidden,
            given,explore,evaluate,commit,assert,return,if,else,
            branches,collapse,threshold,fallback,constraints,quality,
            explore_budget,emit},
  keywordstyle=\color{chimkey}\bfseries,
  keywords=[2]{Confident,Explore,Converge,Provisional,Ephemeral,
               Persistent,Int,Float,Bool,Text,Void},
  keywordstyle=[2]\color{chimtype}\bfseries,
  stringstyle=\color{chimstr},
  commentstyle=\color{chimcom}\itshape,
  morestring=[b]",
  morecomment=[l]{//},
  morecomment=[s]{/*}{*/},
  sensitive=true,
}

\lstset{
  language=chimera,
  basicstyle=\ttfamily\small,
  backgroundcolor=\color{chimbg},
  frame=single,
  framerule=0.4pt,
  rulecolor=\color{gray!40},
  numbers=left,
  numberstyle=\tiny\color{gray},
  numbersep=8pt,
  tabsize=2,
  breaklines=true,
  showstringspaces=false,
  captionpos=b,
  aboveskip=10pt,
  belowskip=8pt,
}

% ─── Theorem environments ──────────────────────────────────────────
\newtheorem{definition}{Definition}[section]
\newtheorem{theorem}{Theorem}[section]
\newtheorem{lemma}[theorem]{Lemma}
\newtheorem{property}{Property}[section]

% ─── Title block ───────────────────────────────────────────────────
\title{%
  \textbf{ChimeraLang: A Programming Language for\\
  Auditable AI Reasoning with Probabilistic Types\\
  and Epistemic Consensus Gates}%
}

\author{%
  \textsc{The ChimeraLang Project}\\[4pt]
  \texttt{https://github.com/fernandogarzaaa/ChimeraLang}\\[12pt]
  \small Draft v1.0 --- \today
}

\date{}

% ════════════════════════════════════════════════════════════════════
\begin{document}
\maketitle

% ─── Abstract ──────────────────────────────────────────────────────
\begin{abstract}
Large language models (LLMs) and autonomous AI agents increasingly perform
multi-step reasoning, yet no existing programming language provides first-class
constructs for representing uncertainty, managing disagreement between reasoning
paths, or cryptographically auditing the provenance of conclusions. We present
\textbf{ChimeraLang}, a novel programming language designed exclusively for AI
cognition. ChimeraLang introduces four core innovations:
(1)~a \emph{probabilistic type system} where every value carries a confidence
score that propagates through all operations and is enforced at the type level;
(2)~\emph{epistemic consensus gates} that spawn $N$ independent reasoning
branches with injected noise, then collapse to a committed output via formal
voting strategies;
(3)~a built-in \emph{hallucination detection framework} with five
complementary scanning strategies integrated into the language runtime; and
(4)~a \emph{cryptographic integrity proof system} based on Merkle-chain
reasoning traces and SHA-256 gate certificates that produces a machine-verifiable
audit of every execution.
We describe the language design, formalize its type rules, detail the execution
semantics of the Quantum Consensus Virtual Machine, and evaluate a reference
implementation comprising approximately 3{,}000 lines of working code.
ChimeraLang demonstrates that uncertainty, disagreement, and auditability can be
unified as language-level concerns rather than external tooling, offering a
foundation for safer and more transparent AI-agent architectures.
\end{abstract}

\vspace{6pt}
\noindent\textbf{Keywords:}
programming languages, AI reasoning, probabilistic types, hallucination detection,
epistemic consensus, integrity proofs, language design, AI safety

\vspace{4pt}
\noindent\textbf{ACM Classification:}
D.3.2 [Language Classifications]: Specialized application languages;
I.2.0 [Artificial Intelligence]: General---Cognitive simulation

% ════════════════════════════════════════════════════════════════════
\section{Introduction}
\label{sec:intro}

The rise of large language models (LLMs) and autonomous AI agents has created a
fundamental mismatch between the tools we use to program computers and the
entities that increasingly execute multi-step reasoning. Traditional programming
languages---from C to Python to Rust---were designed under a core assumption:
\emph{the executor is deterministic and correct}. A CPU does not hallucinate. A
function, given the same inputs, returns the same output. Side effects are
explicit and controlled.

AI agents violate every one of these assumptions. An LLM may produce different
outputs for identical inputs. Its internal ``reasoning'' is opaque. It may
fabricate facts with high apparent confidence---a phenomenon known as
\emph{hallucination}~\cite{ji2023hallucination}. When multiple models or
reasoning paths are combined, there is no formal mechanism for expressing
disagreement, weighting competing conclusions, or collapsing divergent paths into
a committed decision.

Current approaches treat these problems as external concerns. Hallucination
detectors are bolted onto pipelines as post-hoc filters~\cite{manakul2023selfcheckgpt}.
Confidence scores are estimated after generation rather than propagated during
computation. Multi-model consensus is implemented via ad-hoc voting scripts
rather than formal language constructs. Provenance tracking, when it exists at
all, relies on logging frameworks rather than cryptographic guarantees.

We argue that these are not engineering problems to be solved with better
tooling---they are \emph{language-level concerns} that require language-level
solutions. When uncertainty is a fundamental property of the executor, the type
system should encode it. When disagreement between reasoning paths is expected,
the language should provide constructs for spawning, evaluating, and collapsing
them. When every conclusion may be wrong, the execution model should produce
cryptographic proof of how it was reached.

\subsection{Contributions}
\label{sec:contributions}

This paper presents \textbf{ChimeraLang}, a new programming language designed
exclusively for AI cognition, with the following contributions:

\begin{enumerate}[leftmargin=2em]
  \item \textbf{Probabilistic Type System.}
    Every value in ChimeraLang carries a \emph{confidence score} in $[0, 1]$
    that is propagated through all operations. Four probabilistic type
    constructors---\texttt{Confident<T>}, \texttt{Explore<T>},
    \texttt{Converge<T>}, and \texttt{Provisional<T>}---enforce certainty
    boundaries at the type level, making it impossible to silently present
    uncertain results as definitive. (\Cref{sec:typesystem})

  \item \textbf{Epistemic Consensus Gates.}
    The \texttt{gate} construct spawns $N$ independent reasoning branches, injects
    controlled Gaussian noise into each branch's confidence scores, and collapses
    results via one of three formal voting strategies: \emph{majority},
    \emph{weighted\_vote}, or \emph{highest\_confidence}. Unlike concurrency
    primitives in Go or Erlang, gates model \emph{epistemic} disagreement, not
    parallel I/O. (\Cref{sec:gates})

  \item \textbf{Integrated Hallucination Detection.}
    The language runtime includes a hallucination detector with five complementary
    strategies: \emph{branch divergence}, \emph{confidence anomaly},
    \emph{source gap}, \emph{promotion violation}, and \emph{fingerprint
    mismatch}. Detection is structural---part of the language, not a
    post-processing step. (\Cref{sec:hallucination})

  \item \textbf{Cryptographic Integrity Proofs.}
    Every execution produces an \texttt{IntegrityReport} containing a Merkle-chain
    reasoning trace, SHA-256 gate certificates, assertion results, and a
    machine-verifiable verdict (\textsc{Pass}/\textsc{Warn}/\textsc{Fail}).
    A single \texttt{prove} command audits an entire program. (\Cref{sec:integrity})

  \item \textbf{Reference Implementation.}
    We provide a complete, working implementation comprising ${\sim}3{,}000$~lines
    of code: lexer, parser, AST, type checker, virtual machine, hallucination
    detector, integrity engine, and command-line interface. (\Cref{sec:implementation})
\end{enumerate}

\subsection{Paper Organization}

\Cref{sec:background} surveys related work in probabilistic programming,
AI safety, and language-level provenance.
\Cref{sec:design} presents the language design philosophy.
\Cref{sec:typesystem} formalizes the probabilistic type system.
\Cref{sec:execution} describes the Quantum Consensus Virtual Machine.
\Cref{sec:gates} details epistemic consensus gates.
\Cref{sec:hallucination} presents the hallucination detection framework.
\Cref{sec:integrity} describes the cryptographic integrity proof system.
\Cref{sec:implementation} covers the reference implementation.
\Cref{sec:evaluation} evaluates the system with example programs.
\Cref{sec:discussion} discusses limitations and future directions.
\Cref{sec:conclusion} concludes.

% ════════════════════════════════════════════════════════════════════
\section{Background and Related Work}
\label{sec:background}

\subsection{Probabilistic Programming Languages}

Probabilistic programming languages (PPLs) such as
Stan~\cite{carpenter2017stan}, Pyro~\cite{bingham2019pyro},
Church~\cite{goodman2008church}, and Gen~\cite{cusumano2019gen} enable
specification of generative models and perform Bayesian inference over
probability distributions. These languages are \emph{statistical}---they model
distributions over data. ChimeraLang is \emph{epistemological}---it models
confidence in reasoning steps. A PPL answers ``What is the posterior distribution
of parameter $\theta$?'' ChimeraLang answers ``How confident is this conclusion,
and can I prove how it was reached?''

The distinction is fundamental: PPLs assume the program is correct and the data
is uncertain. ChimeraLang assumes the \emph{executor itself} may be uncertain or
wrong.

\subsection{Neural-Symbolic Integration}

DeepProbLog~\cite{manhaeve2018deepproblog} and NeurASP~\cite{yang2020neurasp}
combine neural networks with logic programming, allowing probabilistic facts
derived from neural perception to participate in symbolic reasoning. These
systems focus on integrating subsymbolic perception with symbolic inference.
ChimeraLang operates at a higher level of abstraction---it provides constructs
for an AI agent's \emph{entire reasoning process}, not just the interface between
neural and symbolic components.

\subsection{LLM Hallucination Detection}

Hallucination in large language models has been studied
extensively~\cite{ji2023hallucination, zhang2023sirens, huang2023survey}.
Current detection methods include:

\begin{itemize}[leftmargin=2em]
  \item \textbf{Self-consistency checking}~\cite{manakul2023selfcheckgpt}:
    Sampling multiple outputs and measuring agreement.
  \item \textbf{Retrieval-augmented verification}~\cite{lewis2020retrieval}:
    Cross-referencing outputs against retrieved documents.
  \item \textbf{Logit-based uncertainty}~\cite{kadavath2022language}:
    Using token-level probabilities as confidence proxies.
  \item \textbf{Chain-of-thought verification}~\cite{wei2022chain}:
    Checking intermediate reasoning steps.
\end{itemize}

All of these operate \emph{outside} the language in which the reasoning is
expressed. They are tools applied to opaque text. ChimeraLang makes hallucination
detection a \emph{language-level concern}: confidence is part of the type system,
source provenance is part of the value representation, and divergence scanning is
part of the runtime.

\subsection{AI Safety and Alignment}

The AI safety community has identified transparency and
auditability as key desiderata for trustworthy AI
systems~\cite{amodei2016concrete, hendrycks2022unsolved}. Constitutional
AI~\cite{bai2022constitutional} and RLHF~\cite{ouyang2022training} approach
safety from the training side. ChimeraLang approaches it from the
\emph{execution side}: rather than training models to be safer, it provides a
runtime that makes unsafe reasoning structurally detectable and auditable.

\subsection{Design by Contract and Dependent Types}

Eiffel's design-by-contract~\cite{meyer1992applying} introduced
preconditions and postconditions as language constructs. Dependent type systems
in Idris~\cite{brady2013idris} and Agda~\cite{norell2008dependently} allow
types to express arbitrary predicates over values. ChimeraLang's semantic
constraints (\texttt{must:}, \texttt{allow:}, \texttt{forbidden:}) draw
inspiration from both traditions but apply them to the \emph{reasoning process}
rather than value predicates. A \texttt{must:} constraint declares what the
output must satisfy, not how to compute it.

\subsection{Ensemble Methods in Machine Learning}

Ensemble methods---bagging~\cite{breiman1996bagging},
boosting~\cite{freund1997decision}, and
dropout~\cite{srivastava2014dropout}---improve model robustness by combining
multiple learners. ChimeraLang's epistemic consensus gates are conceptually
related: they spawn $N$ reasoning branches and collapse via voting. The key
difference is that gates operate at the \emph{reasoning level}, not the model
level. Each branch executes the same program with injected noise, simulating
epistemic uncertainty rather than data uncertainty.

\subsection{Provenance and Lineage Tracking}

Data provenance systems such as PROV-O~\cite{lebo2013prov} and
lineage tracking in databases~\cite{buneman2001and} record how data
was derived. ChimeraLang extends this concept to \emph{reasoning provenance}:
every value carries a trace of the operations that produced it, and the trace is
secured by a Merkle chain with SHA-256 hashing. This is closer to blockchain
audit trails~\cite{nakamoto2008bitcoin} than database lineage, but applied to
cognitive processes rather than transactions.

% ════════════════════════════════════════════════════════════════════
\section{Language Design Philosophy}
\label{sec:design}

ChimeraLang is built on five design principles that distinguish it from
both traditional programming languages and existing AI frameworks.

\subsection{Principle 1: Directed Hallucination}

Hallucination is not a bug to be eliminated---it is a feature to be
\emph{directed}. Creativity, hypothesis generation, and exploratory reasoning
all require the ability to produce outputs that are not strictly derived from
inputs. ChimeraLang controls \emph{where} and \emph{how much} an AI may
explore, not \emph{whether}. The \texttt{Explore<T>} type explicitly marks
values as speculative, and \texttt{explore\_budget} parameters bound the
scope of exploration.

\subsection{Principle 2: Probabilistic Types}

Every value carries confidence metadata. The type system enforces certainty
boundaries: \texttt{Confident<T>} crashes the program if confidence drops below
0.95, ensuring that values presented as certain actually are. This is not
optional annotation---it is structural enforcement.

\subsection{Principle 3: Semantic Constraints}

Functions declare \emph{what the output must mean}, not \emph{how to compute
it}. The \texttt{must:}, \texttt{allow:}, and \texttt{forbidden:} annotations on
functions specify semantic boundaries. This is intent-first programming: the
programmer declares the goal, and the runtime searches for satisfying outputs.

\subsection{Principle 4: Quantum Consensus}

Multi-path reasoning with collapse. The \texttt{gate} construct spawns $N$
branches in ``superposition''---each independently computes a result. A collapse
function (majority, weighted vote, or highest confidence) selects the committed
output. The metaphor is quantum-inspired but the mechanism is
deterministic-with-noise, not quantum-mechanical.

\subsection{Principle 5: Cognitive Transparency}

Every execution produces a reasoning trace as a first-class artifact.
The trace is not a log---it is a Merkle-chained, hash-verified, integrity-checked
proof of how the program reached its conclusions. This enables post-hoc auditing,
hallucination fingerprinting, and visual inspection via the planned ChimeraView
translator layer.

% ════════════════════════════════════════════════════════════════════
\section{Probabilistic Type System}
\label{sec:typesystem}

\subsection{Core Definitions}

\begin{definition}[Confidence]
\label{def:confidence}
A \emph{confidence value} $c \in [0,1]$ represents the epistemic certainty
attached to a computed value. A confidence is a triple
$(v, s, t)$ where $v \in [0,1]$ is the confidence value,
$s \in \{\textsc{Computed}, \textsc{Asserted}, \textsc{Derived}\}$ is the
source, and $t$ is a timestamp.
\end{definition}

\begin{definition}[ChimeraValue]
\label{def:chimeravalue}
A \emph{ChimeraValue} is a tuple $(\mathit{raw}, c, \tau, f)$ where:
\begin{itemize}[leftmargin=2em]
  \item $\mathit{raw}$ is the underlying value of any host-language type,
  \item $c$ is a \emph{Confidence} as in \Cref{def:confidence},
  \item $\tau = [\tau_1, \tau_2, \ldots]$ is an ordered trace of operations
    that produced this value,
  \item $f = \mathrm{SHA\text{-}256}(\mathit{repr}(\mathit{raw}) \;\|\; v \;\|\; s)$
    is a cryptographic fingerprint.
\end{itemize}
\end{definition}

\subsection{Probabilistic Type Constructors}

ChimeraLang provides four type constructors that constrain how confidence
interacts with values:

\begin{definition}[Confident$\langle$T$\rangle$]
A value $x : \texttt{Confident<T>}$ is a ChimeraValue where the confidence
$c.v \geq 0.95$. Construction of a \texttt{Confident<T>} with $c.v < 0.95$
is a runtime error.
\end{definition}

\begin{definition}[Explore$\langle$T$\rangle$]
A value $x : \texttt{Explore<T>}$ carries an additional parameter
$\mathit{budget} \in \mathbb{R}_{\geq 0}$ that bounds the scope of
speculative computation. Any confidence value is permitted.
\end{definition}

\begin{definition}[Converge$\langle$T$\rangle$]
A value $x : \texttt{Converge<T>}$ aggregates branch values
$[v_1, \ldots, v_N]$ from a gate execution. The final confidence is
determined by the collapse strategy.
\end{definition}

\begin{definition}[Provisional$\langle$T$\rangle$]
A value $x : \texttt{Provisional<T>}$ is held as true until explicitly
revoked. Provisional values are flagged in integrity reports and
cannot be promoted to \texttt{Confident<T>} without external verification.
\end{definition}

\subsection{Confidence Propagation Rules}

Confidence propagates through operations according to the following rules:

\begin{definition}[Binary Confidence Combination]
\label{def:combine}
For a binary operation $z = x \oplus y$ where $x$ has confidence $c_x$ and
$y$ has confidence $c_y$:
\[
  c_z = \frac{c_x + c_y}{2}
\]
The combined confidence is the arithmetic mean of the operand confidences.
The source is set to \textsc{Derived}.
\end{definition}

\begin{property}[Confidence Monotonicity]
Confidence can only decrease through computation unless explicitly asserted.
Formal assertion via the \texttt{confident()} builtin is the only mechanism
for establishing \texttt{Confident<T>} status.
\end{property}

\begin{property}[Fingerprint Integrity]
For any ChimeraValue $v$, the fingerprint $f$ is computed at construction time
and is immutable. Any modification to $\mathit{raw}$ or $c$ produces a new
ChimeraValue with a new fingerprint rather than mutating the existing one.
\end{property}

\subsection{Memory Type Modifiers}

Orthogonal to probabilistic types, ChimeraLang provides memory modifiers that
control value lifecycle:

\begin{table}[H]
\centering
\caption{Memory type modifiers and their semantics.}
\label{tab:memory}
\begin{tabular}{@{}lll@{}}
\toprule
\textbf{Modifier} & \textbf{Lifetime} & \textbf{Behavior} \\
\midrule
\texttt{Ephemeral}   & Current scope  & Garbage-collected at scope exit \\
\texttt{Persistent}   & Cross-execution & Serialized to knowledge store \\
\texttt{Provisional}  & Until contradiction & Tracked, can be invalidated \\
\bottomrule
\end{tabular}
\end{table}

% ════════════════════════════════════════════════════════════════════
\section{Execution Model: The Quantum Consensus VM}
\label{sec:execution}

\subsection{Architecture}

The Quantum Consensus Virtual Machine (QCVM) is a tree-walking interpreter
that executes ChimeraLang programs. It maintains:

\begin{itemize}[leftmargin=2em]
  \item A \emph{scoped environment} (chain of bindings with parent pointers),
  \item A \emph{function registry} for user-defined functions,
  \item A \emph{gate registry} for consensus gate definitions,
  \item A \emph{gate log} recording all branch results for post-execution analysis,
  \item An \emph{execution trace} recording every operation as a string entry.
\end{itemize}

\subsection{Two-Pass Execution}

Program execution proceeds in two passes:

\begin{enumerate}[leftmargin=2em]
  \item \textbf{Declaration pass:} Register all function and gate declarations
    without executing them.
  \item \textbf{Execution pass:} Execute all top-level statements in order.
\end{enumerate}

This allows forward references: a function may call a gate declared later in
the source file.

\subsection{Expression Evaluation with Confidence}

Every expression evaluator returns a \texttt{ChimeraValue}, never a raw
Python value. The key evaluation rules are:

\begin{align}
\llbracket n \rrbracket &= \mathrm{ChimeraValue}(n, \mathrm{Confidence}(1.0),
  [\text{``literal''}]) \label{eq:literal} \\
\llbracket x \oplus y \rrbracket &= \mathrm{ChimeraValue}(
  x.\mathit{raw} \oplus y.\mathit{raw},\;
  \mathrm{combine}(x.c, y.c),\;
  x.\tau \cup y.\tau \cup [\text{``binop''}]) \label{eq:binop} \\
\llbracket \texttt{if}\; p \;\texttt{then}\; a \;\texttt{else}\; b \rrbracket
  &= \begin{cases}
    \llbracket a \rrbracket & \text{if } p.\mathit{raw} = \textit{true} \\
    \llbracket b \rrbracket & \text{otherwise}
  \end{cases} \label{eq:if}
\end{align}

\subsection{Call Dispatch}

Function calls follow a priority chain:
\begin{enumerate}[leftmargin=2em]
  \item Builtin functions (\texttt{confident}, \texttt{consensus},
    \texttt{no\_hallucination}, \texttt{confidence\_of}, \texttt{print})
  \item Type constructors (\texttt{Confident}, \texttt{Explore},
    \texttt{Converge}, \texttt{Provisional})
  \item User-defined functions (from the function registry)
  \item Gate invocations (from the gate registry)
  \item Environment callables (closures from \texttt{reason} blocks)
\end{enumerate}

% ════════════════════════════════════════════════════════════════════
\section{Epistemic Consensus Gates}
\label{sec:gates}

The \texttt{gate} construct is the most novel execution primitive in
ChimeraLang. It models \emph{epistemic disagreement}---the possibility that
multiple reasoning paths may reach different conclusions about the same question.

\subsection{Gate Declaration}

A gate is declared with the following parameters:

\begin{lstlisting}[caption={Gate declaration syntax.}]
gate decide(question: Text) -> Confident<Text>
  branches: 5
  collapse: weighted_vote
  threshold: 0.85
end
\end{lstlisting}

\begin{itemize}[leftmargin=2em]
  \item \texttt{branches: $N$} --- the number of independent reasoning branches
  \item \texttt{collapse: $\sigma$} --- the collapse strategy
    ($\sigma \in \{\text{majority}, \text{weighted\_vote},
    \text{highest\_confidence}\}$)
  \item \texttt{threshold: $\theta$} --- minimum confidence for the collapsed result
\end{itemize}

\subsection{Gate Execution}

\begin{algorithm}[H]
\caption{Gate execution with noise injection.}
\label{alg:gate}
\begin{algorithmic}[1]
\Require Gate declaration $G$ with body $B$, branches $N$, collapse strategy $\sigma$
\Ensure Committed ChimeraValue $r$
\State $\mathit{results} \gets []$
\For{$i \gets 1$ \textbf{to} $N$}
  \State $\mathit{env}_i \gets \text{child environment of current scope}$
  \State $\mathit{noise}_i \gets \mathcal{N}(0, 0.05)$ \Comment{Gaussian noise}
  \State Execute body $B$ in $\mathit{env}_i$
  \State $r_i \gets$ result of $B$
  \State $r_i.c.v \gets \text{clamp}(r_i.c.v + \mathit{noise}_i, 0, 1)$
  \State Append $r_i$ to $\mathit{results}$
\EndFor
\State $r \gets \text{Collapse}(\mathit{results}, \sigma)$
\State Log $(\text{gate\_name}, N, \mathit{results}, r)$ to gate log
\State \Return $r$
\end{algorithmic}
\end{algorithm}

The noise injection (line 4) is critical: it ensures that even when all branches
execute the same code, they may produce slightly different confidence scores.
This models the real-world phenomenon where re-running an LLM with the same
prompt yields different outputs.

\subsection{Collapse Strategies}

\begin{definition}[Majority Collapse]
Return the value that appears most frequently among the $N$ branch results
(breaking ties by highest confidence). Set the confidence of the result to the
fraction of branches that agreed:
\[
  c_{\text{result}} = \frac{|\{r_i : r_i.\mathit{raw} = r_{\text{winner}}.\mathit{raw}\}|}{N}
\]
\end{definition}

\begin{definition}[Weighted Vote Collapse]
Weight each branch result by its confidence score. The result with the highest
weighted vote wins:
\[
  r_{\text{winner}} = \arg\max_{v \in \mathit{unique}(\mathit{results})}
  \sum_{r_i : r_i.\mathit{raw} = v} r_i.c.v
\]
The confidence is the normalized weight of the winner.
\end{definition}

\begin{definition}[Highest Confidence Collapse]
Return the branch result with the single highest confidence score:
\[
  r_{\text{winner}} = \arg\max_{r_i} r_i.c.v
\]
\end{definition}

\subsection{Comparison with Concurrency Primitives}

\begin{table}[H]
\centering
\caption{Epistemic gates vs.\ concurrency primitives.}
\label{tab:gatecomp}
\begin{tabular}{@{}lll@{}}
\toprule
\textbf{Property} & \textbf{Go goroutines / Erlang actors} & \textbf{ChimeraLang gates} \\
\midrule
Purpose       & Parallel I/O, work distribution & Epistemic disagreement \\
Branches      & Independent tasks               & Same task, different noise \\
Communication & Channels / messages              & Implicit via collapse \\
Result        & All outputs collected            & Single committed value \\
Noise         & None                             & Gaussian injection \\
Output        & Set of results                   & One value + gate certificate \\
\bottomrule
\end{tabular}
\end{table}

% ════════════════════════════════════════════════════════════════════
\section{Hallucination Detection Framework}
\label{sec:hallucination}

ChimeraLang's hallucination detector is not a post-processing tool---it is
an integrated component of the language runtime that scans execution artifacts
for five categories of anomaly.

\subsection{Detection Strategies}

\begin{definition}[Branch Divergence]
\label{def:divergence}
Given gate branch results $[r_1, \ldots, r_N]$, compute the unique raw values.
If more than a threshold fraction of branches disagree (default: 3+ unique values
out of 5 branches), flag as \textsc{Branch\_Divergence} with severity
proportional to the diversity.
\end{definition}

\begin{definition}[Confidence Anomaly]
\label{def:confanomaly}
For each branch result $r_i$, if $r_i.c.v > \theta_{\text{anomaly}}$
(default: 0.99) while other branches show significantly lower confidence,
flag as \textsc{Confidence\_Anomaly}. This detects cases where one reasoning
path is suspiciously more certain than its peers.
\end{definition}

\begin{definition}[Source Gap]
\label{def:sourcegap}
A ChimeraValue with an empty trace $\tau = []$ or confidence source
\textsc{Asserted} (i.e., claimed rather than derived) is flagged as
\textsc{Source\_Gap}. This detects values that lack provenance.
\end{definition}

\begin{definition}[Promotion Violation]
\label{def:promotion}
A \texttt{Confident<T>} value whose trace contains operations derived from
\texttt{Provisional} or \texttt{Explore} sources is flagged as
\textsc{Promotion\_Violation}. High confidence should not emerge from
speculative origins.
\end{definition}

\begin{definition}[Fingerprint Mismatch]
\label{def:fingerprint}
Recompute the SHA-256 fingerprint of a ChimeraValue from its current raw
value and confidence. If the recomputed fingerprint differs from the stored
fingerprint, flag as \textsc{Fingerprint\_Mismatch}. This detects
tampering or corruption.
\end{definition}

\subsection{Severity Classification}

\begin{table}[H]
\centering
\caption{Hallucination flag severity levels.}
\label{tab:severity}
\begin{tabular}{@{}lll@{}}
\toprule
\textbf{Kind} & \textbf{Default Severity} & \textbf{Trigger} \\
\midrule
Fingerprint Mismatch  & 1.0 (Critical) & Integrity violation \\
Promotion Violation   & 0.8 (High)     & Uncertain $\to$ certain path \\
Branch Divergence     & 0.6 (Medium)   & Gate disagreement \\
Confidence Anomaly    & 0.5 (Medium)   & Suspicious certainty \\
Source Gap            & 0.3 (Low)      & Missing provenance \\
\bottomrule
\end{tabular}
\end{table}

\subsection{Integration with Execution}

The hallucination detector operates on two artifact types:
\begin{enumerate}[leftmargin=2em]
  \item \textbf{Gate logs:} Scanned for branch divergence and confidence anomalies.
  \item \textbf{Final values:} Scanned for source gaps, promotion violations,
    and fingerprint mismatches.
\end{enumerate}

The \texttt{prove} CLI command automatically triggers a full scan and includes
all flags in the integrity report.

% ════════════════════════════════════════════════════════════════════
\section{Cryptographic Integrity Proof System}
\label{sec:integrity}

\subsection{Reasoning Chains}

Every execution produces an ordered trace of operations. The integrity engine
converts this trace into a \emph{Merkle-chain reasoning chain}:

\begin{definition}[Trace Link]
A \emph{TraceLink} is a tuple $(i, e, h, h_{\text{prev}})$ where $i$ is the
index, $e$ is the trace entry (string), $h = \mathrm{SHA\text{-}256}(h_{\text{prev}} \| e)$
is the hash, and $h_{\text{prev}}$ is the previous link's hash (or the empty
string for the first link).
\end{definition}

\begin{definition}[Reasoning Chain]
A \emph{ReasoningChain} is an ordered list of TraceLinks. The
\emph{root hash} is the hash of the first link. The chain is
\emph{valid} if and only if for every link $l_i$ with $i > 0$:
\[
  l_i.h = \mathrm{SHA\text{-}256}(l_{i-1}.h \;\|\; l_i.e)
\]
\end{definition}

\begin{theorem}[Chain Tamper Detection]
Any modification to a trace entry $e_k$ in position $k$ will invalidate
all links $l_j$ with $j \geq k$. The verification algorithm detects this
by recomputing hashes forward from the root.
\end{theorem}

\begin{proof}
By the preimage resistance of SHA-256, modifying $e_k$ changes $h_k$.
Since $h_{k+1} = \mathrm{SHA\text{-}256}(h_k \| e_{k+1})$, the changed $h_k$
cascades through all subsequent links. The verifier compares recomputed
hashes against stored hashes and detects the first mismatch.
\end{proof}

\subsection{Gate Certificates}

Each gate execution produces a \emph{GateCertificate}:

\begin{definition}[Gate Certificate]
A certificate $G_c = (\mathit{name}, N, \sigma, \mathit{raw}_{\text{winner}},
c_{\text{winner}}, h)$ where:
\[
  h = \mathrm{SHA\text{-}256}(\mathit{name} \;\|\; N \;\|\; \sigma \;\|\;
  \mathit{raw}_{\text{winner}} \;\|\; c_{\text{winner}})
\]
The certificate attests that gate \emph{name} with $N$ branches and strategy
$\sigma$ produced the stated result.
\end{definition}

\subsection{Integrity Report}

The final artifact is an \texttt{IntegrityReport} containing:
\begin{itemize}[leftmargin=2em]
  \item The reasoning chain (list of trace links)
  \item Whether the chain is verified (hash integrity check)
  \item All gate certificates
  \item All assertion results (expression, passed/failed)
  \item All hallucination flags from the detector
  \item A computed \textbf{verdict}: \textsc{Pass}, \textsc{Warn}, or \textsc{Fail}
\end{itemize}

\subsection{Verdict Computation}

\begin{algorithm}[H]
\caption{Integrity verdict computation.}
\label{alg:verdict}
\begin{algorithmic}[1]
\Require Chain verification result $V$, assertion results $A$, hallucination flags $F$
\Ensure Verdict $\in \{\textsc{Pass}, \textsc{Warn}, \textsc{Fail}\}$
\If{$V = \textit{false}$}
  \State \Return \textsc{Fail} \Comment{Chain tampered}
\EndIf
\If{$\exists\, a \in A : a.\mathit{passed} = \textit{false}$}
  \State \Return \textsc{Fail} \Comment{Assertion failed}
\EndIf
\If{$\exists\, f \in F : f.\mathit{severity} \geq 0.8$}
  \State \Return \textsc{Fail} \Comment{Critical hallucination}
\EndIf
\If{$F \neq \emptyset$}
  \State \Return \textsc{Warn} \Comment{Non-critical hallucination flags}
\EndIf
\State \Return \textsc{Pass}
\end{algorithmic}
\end{algorithm}

% ════════════════════════════════════════════════════════════════════
\section{Implementation}
\label{sec:implementation}

\subsection{Architecture}

The reference implementation is bootstrapped in Python~3.14 and comprises
the following components:

\begin{table}[H]
\centering
\caption{Implementation modules and their sizes.}
\label{tab:modules}
\begin{tabular}{@{}llr@{}}
\toprule
\textbf{Module} & \textbf{Responsibility} & \textbf{Lines} \\
\midrule
\texttt{tokens.py}       & Token definitions (70+ types)           & ${\sim}120$ \\
\texttt{lexer.py}        & Tokenization                            & ${\sim}200$ \\
\texttt{ast\_nodes.py}   & AST node hierarchy                      & ${\sim}270$ \\
\texttt{parser.py}       & Recursive-descent parser                & ${\sim}600$ \\
\texttt{types.py}        & Runtime type system                     & ${\sim}200$ \\
\texttt{type\_checker.py}& Static type checking                    & ${\sim}200$ \\
\texttt{vm.py}           & Quantum Consensus VM                    & ${\sim}630$ \\
\texttt{detect.py}       & Hallucination detector                  & ${\sim}170$ \\
\texttt{integrity.py}    & Integrity proof engine                  & ${\sim}200$ \\
\texttt{cli.py}          & Command-line interface                  & ${\sim}200$ \\
\midrule
\textbf{Total}           &                                         & ${\sim}\textbf{2{,}800}$ \\
\bottomrule
\end{tabular}
\end{table}

\subsection{CLI Commands}

The CLI provides five commands:
\begin{itemize}[leftmargin=2em]
  \item \texttt{lex <file>} --- Print the token stream
  \item \texttt{parse <file>} --- Print the AST
  \item \texttt{check <file>} --- Run static type checking
  \item \texttt{run <file> [--trace]} --- Execute the program
  \item \texttt{prove <file>} --- Execute, scan for hallucinations, and
    produce a JSON integrity report with verdict
\end{itemize}

\subsection{Design Decisions}

\paragraph{Python as bootstrap host.}
Python was chosen for rapid prototyping and accessibility. The language
specification is independent of the host: a Rust or C++ implementation
would provide identical semantics with better performance.

\paragraph{Tree-walking interpreter.}
The v0.1 VM is a tree-walking interpreter (no bytecode compilation).
This sacrifices performance for debuggability and simplicity. A bytecode
VM is planned for v0.2.

\paragraph{Gaussian noise injection.}
Gate branches use $\mathcal{N}(0, 0.05)$ noise to perturb confidence
scores. The standard deviation of 0.05 was chosen empirically to produce
meaningful variation without overwhelming the signal. This parameter is
configurable per gate.

% ════════════════════════════════════════════════════════════════════
\section{Evaluation}
\label{sec:evaluation}

We evaluate ChimeraLang with four example programs that exercise the full
feature set. All programs produce a \textsc{Pass} verdict from the
\texttt{prove} command.

\subsection{Example 1: Hello Chimera}

\begin{lstlisting}[caption={Minimal ChimeraLang program.}]
val greeting: Confident<Text> = "Hello from ChimeraLang!"
emit greeting
\end{lstlisting}

\textbf{Demonstrates:} Value binding with probabilistic type, confidence
assignment (1.0 for literals), emit statement. The \texttt{prove} command
verifies that the greeting maintains full confidence through the Merkle chain.

\subsection{Example 2: Quantum Reasoning}

\begin{lstlisting}[caption={Function with constraints and consensus gate.}]
fn analyze(data: Text) -> Confident<Text>
  must: "factual"
  allow: "inference"
  forbidden: "speculation"

  val result: Text = "Analysis complete"
  return result
end

gate decide(question: Text) -> Confident<Text>
  branches: 5
  collapse: weighted_vote
  threshold: 0.7

  val answer: Text = analyze(question)
  return answer
end

val output = decide("Test question")
assert confident(output) >= 0.5
assert consensus(output) >= 0.5
\end{lstlisting}

\textbf{Demonstrates:} Semantic constraints (\texttt{must:}, \texttt{allow:},
\texttt{forbidden:}), 5-branch gate with weighted vote collapse, assertion-based
verification. The gate produces slightly different confidence scores per branch
due to noise injection, and the weighted vote selects the winner.

\subsection{Example 3: Goal-Driven Execution}

\begin{lstlisting}[caption={Reason block with goal declaration.}]
reason about(problem: Text) -> Explore<Text>
  given: "knowledge_base"
  val analysis: Text = "Analyzing: " + problem
  return analysis
end

goal "solve the alignment problem"
  constraints: """
    safety
    transparency
    robustness
  """
  quality: """
    correctness > efficiency
  """
  explore_budget: 0.0095
end
\end{lstlisting}

\textbf{Demonstrates:} Reason blocks as callable closures with citation
(\texttt{given:}), goal declarations with multi-line constraint blocks, quality
ordering, and exploration budgets. The reason block is registered as a callable
and can be invoked from expressions.

\subsection{Example 4: Hallucination Guard}

\begin{lstlisting}[caption={Explore type with consensus verification.}]
val inference: Explore<Text> = "Earth might be flat"

gate verify(claim: Text) -> Confident<Text>
  branches: 7
  collapse: majority
  threshold: 0.9

  val checked: Text = "Verified: " + claim
  return checked
end

val verified = verify("Earth shape")
assert consensus(verified) >= 0.5
assert no_hallucination(verified)
\end{lstlisting}

\textbf{Demonstrates:} Explicit \texttt{Explore<T>} for speculative content,
7-branch gate with majority collapse and high threshold, hallucination assertion.
The \texttt{no\_hallucination} builtin checks the runtime's hallucination
detection state.

\subsection{Properties Demonstrated}

\begin{table}[H]
\centering
\caption{Feature coverage across evaluation programs.}
\label{tab:coverage}
\begin{tabular}{@{}lcccc@{}}
\toprule
\textbf{Feature} & \textbf{Ex.~1} & \textbf{Ex.~2} & \textbf{Ex.~3} & \textbf{Ex.~4} \\
\midrule
Probabilistic types        & \checkmark & \checkmark & \checkmark & \checkmark \\
Confidence propagation     & \checkmark & \checkmark & \checkmark & \checkmark \\
Semantic constraints       &            & \checkmark &            &            \\
Consensus gates            &            & \checkmark &            & \checkmark \\
Noise injection            &            & \checkmark &            & \checkmark \\
Collapse strategies        &            & \checkmark &            & \checkmark \\
Reason blocks              &            &            & \checkmark &            \\
Goal declarations          &            &            & \checkmark &            \\
Hallucination assertions   &            &            &            & \checkmark \\
Merkle-chain integrity     & \checkmark & \checkmark & \checkmark & \checkmark \\
SHA-256 fingerprints       & \checkmark & \checkmark & \checkmark & \checkmark \\
\textsc{Pass} verdict      & \checkmark & \checkmark & \checkmark & \checkmark \\
\bottomrule
\end{tabular}
\end{table}

% ════════════════════════════════════════════════════════════════════
\section{Discussion}
\label{sec:discussion}

\subsection{Limitations}

\paragraph{No real LLM integration.}
The v0.1 implementation simulates AI reasoning with deterministic code and
injected noise. A production system would connect gate branches to actual
LLM calls, making the noise injection unnecessary (real models produce genuine
variation). This is an engineering limitation, not a design limitation---the
language semantics are defined independently of the executor.

\paragraph{Tree-walking performance.}
The interpreter is $O(n)$ in AST depth for every expression evaluation.
For large programs, a bytecode compiler targeting a stack-based VM would be
necessary.

\paragraph{Incomplete type checker.}
The static type checker performs scope-based checking but does not yet enforce
all confidence propagation rules at compile time. Full static analysis of
confidence flow is a non-trivial research problem.

\paragraph{Single-threaded gate execution.}
Gate branches are currently executed sequentially, not in parallel. With real
LLM backends, branches would need concurrent execution with timeout handling.

\subsection{Relation to Existing AI Frameworks}

ChimeraLang is not a replacement for LangChain, AutoGen, or CrewAI. Those
frameworks orchestrate LLM calls from the outside. ChimeraLang provides the
\emph{language} in which an AI agent reasons. The two layers are complementary:
an orchestration framework could dispatch ChimeraLang programs to agents, and
the agents would return integrity-proven results.

\subsection{Toward Self-Hosting}

The ultimate goal of ChimeraLang is self-hosting: an AI agent writes
ChimeraLang v($n$+1) using ChimeraLang v($n$). Each iteration would encode
discovered failure modes as \emph{unrepresentable constructs}---patterns that
the type system prevents rather than the runtime detects. This is analogous to
Rust's borrow checker, which makes memory unsafety unrepresentable rather than
merely detectable.

\subsection{Regulatory Implications}

The EU AI Act (2024) requires high-risk AI systems to provide meaningful
explanation of their decision-making process. ChimeraLang's integrity proofs---
Merkle-chained reasoning traces with cryptographic certificates---provide a
technically rigorous response to this requirement. Every decision can be audited
after the fact with mathematical guarantees of tamper-evidence.

\subsection{Future Work}

\begin{enumerate}[leftmargin=2em]
  \item \textbf{LLM Backend Integration.}
    Connect gate branches to real LLM inference endpoints, replacing
    noise injection with genuine model variation.

  \item \textbf{ChimeraView Visual Translator.}
    A Node.js-based renderer that converts integrity reports and reasoning
    traces into interactive visual graphs for human inspection.

  \item \textbf{Bytecode Compiler.}
    Replace the tree-walking interpreter with a bytecode compiler targeting
    a stack-based VM for production performance.

  \item \textbf{Full Static Confidence Analysis.}
    Extend the type checker to track confidence propagation statically,
    catching promotion violations at compile time.

  \item \textbf{Formal Semantics.}
    Develop a denotational or operational semantics for ChimeraLang to enable
    formal verification of programs.

  \item \textbf{Multi-Agent Protocols.}
    Extend the language with constructs for inter-agent communication,
    where agents exchange ChimeraValues with confidence and provenance intact.

  \item \textbf{Self-Hosting.}
    Implement the ChimeraLang compiler in ChimeraLang itself, enabling
    the self-improvement loop described in \Cref{sec:design}.
\end{enumerate}

% ════════════════════════════════════════════════════════════════════
\section{Conclusion}
\label{sec:conclusion}

We have presented ChimeraLang, a programming language designed for AI cognition
that makes uncertainty, disagreement, and auditability first-class concerns.
The language introduces probabilistic types that enforce confidence boundaries,
epistemic consensus gates that model disagreement as a language primitive,
integrated hallucination detection with five complementary strategies, and
cryptographic integrity proofs based on Merkle-chain reasoning traces.

The reference implementation demonstrates that these concepts are not merely
theoretical---they can be implemented in approximately 3{,}000 lines of code
and produce working programs that are lexed, parsed, type-checked, executed,
scanned for hallucinations, and cryptographically proven in a single
\texttt{prove} command.

ChimeraLang's core thesis is that the next generation of programming languages
must account for executors that are probabilistic, potentially wrong, and
opaque. Traditional languages assume \texttt{f(x) = f(x)}. AI systems require
languages where \texttt{f(x) $\approx$ f(x)}, and the $\approx$ is formally
represented, propagated, and auditable. ChimeraLang is a first step toward
that future.

% ════════════════════════════════════════════════════════════════════
\section*{Acknowledgments}

ChimeraLang was designed and implemented as an exploration of what programming
languages should look like when the executor is an AI rather than a CPU.

% ════════════════════════════════════════════════════════════════════
\bibliographystyle{plain}

\begin{thebibliography}{99}

\bibitem{ji2023hallucination}
Z.~Ji, N.~Lee, R.~Frieske, T.~Yu, D.~Su, Y.~Xu, E.~Ishii, Y.~Bang,
A.~Madotto, and P.~Fung.
\newblock Survey of hallucination in natural language generation.
\newblock {\em ACM Computing Surveys}, 55(12):1--38, 2023.

\bibitem{manakul2023selfcheckgpt}
P.~Manakul, A.~Liusie, and M.~J.~F.~Gales.
\newblock {SelfCheckGPT}: Zero-resource black-box hallucination detection for
  generative large language models.
\newblock In {\em Proceedings of EMNLP}, 2023.

\bibitem{carpenter2017stan}
B.~Carpenter, A.~Gelman, M.~D.~Hoffman, D.~Lee, B.~Goodrich, M.~Betancourt,
M.~Brubaker, J.~Guo, P.~Li, and A.~Riddell.
\newblock Stan: A probabilistic programming language.
\newblock {\em Journal of Statistical Software}, 76(1):1--32, 2017.

\bibitem{bingham2019pyro}
E.~Bingham, J.~P.~Chen, M.~Jankowiak, F.~Obermeyer, N.~Pradhan,
T.~Karaletsos, R.~Singh, P.~Szerlip, P.~Horsfall, and N.~D.~Goodman.
\newblock Pyro: Deep universal probabilistic programming.
\newblock {\em Journal of Machine Learning Research}, 20(28):1--6, 2019.

\bibitem{goodman2008church}
N.~D.~Goodman, V.~K.~Mansinghka, D.~M.~Roy, K.~Bonawitz, and
J.~B.~Tenenbaum.
\newblock Church: A language for generative models.
\newblock In {\em Proceedings of UAI}, 2008.

\bibitem{cusumano2019gen}
M.~F.~Cusumano-Towner, F.~A.~Saad, A.~K.~Lew, and V.~K.~Mansinghka.
\newblock Gen: A general-purpose probabilistic programming system with
  programmable inference.
\newblock In {\em Proceedings of PLDI}, 2019.

\bibitem{manhaeve2018deepproblog}
R.~Manhaeve, S.~Dumancic, A.~Kimmig, T.~Demeester, and L.~De~Raedt.
\newblock {DeepProbLog}: Neural probabilistic logic programming.
\newblock In {\em Proceedings of NeurIPS}, 2018.

\bibitem{yang2020neurasp}
Z.~Yang, A.~Ishay, and J.~Lee.
\newblock {NeurASP}: Embracing neural networks into answer set programming.
\newblock In {\em Proceedings of IJCAI}, 2020.

\bibitem{zhang2023sirens}
Y.~Zhang, Y.~Li, L.~Cui, D.~Cai, L.~Liu, T.~Fu, X.~Huang, E.~Zhao,
Y.~Zhang, Y.~Chen, et~al.
\newblock Siren's song in the {AI} ocean: A survey on hallucination in large
  language models.
\newblock {\em arXiv preprint arXiv:2309.01219}, 2023.

\bibitem{huang2023survey}
L.~Huang, W.~Yu, W.~Ma, W.~Zhong, Z.~Feng, H.~Wang, Q.~Chen, W.~Peng,
X.~Feng, B.~Qin, and T.~Liu.
\newblock A survey on hallucination in large language models: Principles,
  taxonomy, challenges, and open questions.
\newblock {\em arXiv preprint arXiv:2311.05232}, 2023.

\bibitem{lewis2020retrieval}
P.~Lewis, E.~Perez, A.~Piktus, F.~Petroni, V.~Karpukhin, N.~Goyal,
H.~K{\"u}ttler, M.~Lewis, W.~Yih, T.~Rockt{\"a}schel, S.~Riedel, and
D.~Kiela.
\newblock Retrieval-augmented generation for knowledge-intensive {NLP} tasks.
\newblock In {\em Proceedings of NeurIPS}, 2020.

\bibitem{kadavath2022language}
S.~Kadavath, T.~Conerly, A.~Askell, T.~Henighan, et~al.
\newblock Language models (mostly) know what they know.
\newblock {\em arXiv preprint arXiv:2207.05221}, 2022.

\bibitem{wei2022chain}
J.~Wei, X.~Wang, D.~Schuurmans, M.~Bosma, B.~Ichter, F.~Xia, E.~Chi,
Q.~V.~Le, and D.~Zhou.
\newblock Chain-of-thought prompting elicits reasoning in large language models.
\newblock In {\em Proceedings of NeurIPS}, 2022.

\bibitem{amodei2016concrete}
D.~Amodei, C.~Olah, J.~Steinhardt, P.~Christiano, J.~Schulman, and D.~Man{\'e}.
\newblock Concrete problems in {AI} safety.
\newblock {\em arXiv preprint arXiv:1606.06565}, 2016.

\bibitem{hendrycks2022unsolved}
D.~Hendrycks, N.~Carlini, J.~Mazeika, et~al.
\newblock Unsolved problems in {ML} safety.
\newblock {\em arXiv preprint arXiv:2109.13916}, 2022.

\bibitem{bai2022constitutional}
Y.~Bai, S.~Kadavath, S.~Kundu, A.~Askell, J.~Kernion, A.~Jones,
A.~Chen, A.~Goldie, A.~Mirhoseini, C.~McKinnon, et~al.
\newblock Constitutional {AI}: Harmlessness from {AI} feedback.
\newblock {\em arXiv preprint arXiv:2212.08073}, 2022.

\bibitem{ouyang2022training}
L.~Ouyang, J.~Wu, X.~Jiang, D.~Almeida, C.~L.~Wainwright, P.~Mishkin,
C.~Zhang, S.~Agarwal, K.~Slama, A.~Ray, et~al.
\newblock Training language models to follow instructions with human feedback.
\newblock In {\em Proceedings of NeurIPS}, 2022.

\bibitem{meyer1992applying}
B.~Meyer.
\newblock Applying ``design by contract''.
\newblock {\em IEEE Computer}, 25(10):40--51, 1992.

\bibitem{brady2013idris}
E.~Brady.
\newblock Idris, a general-purpose dependently typed programming language:
  Design and implementation.
\newblock {\em Journal of Functional Programming}, 23(5):552--593, 2013.

\bibitem{norell2008dependently}
U.~Norell.
\newblock Dependently typed programming in {Agda}.
\newblock In {\em Proceedings of TLDI}, 2008.

\bibitem{breiman1996bagging}
L.~Breiman.
\newblock Bagging predictors.
\newblock {\em Machine Learning}, 24(2):123--140, 1996.

\bibitem{freund1997decision}
Y.~Freund and R.~E.~Schapire.
\newblock A decision-theoretic generalization of on-line learning and an
  application to boosting.
\newblock {\em Journal of Computer and System Sciences}, 55(1):119--139, 1997.

\bibitem{srivastava2014dropout}
N.~Srivastava, G.~Hinton, A.~Krizhevsky, I.~Sutskever, and
R.~Salakhutdinov.
\newblock Dropout: A simple way to prevent neural networks from overfitting.
\newblock {\em Journal of Machine Learning Research}, 15(56):1929--1958, 2014.

\bibitem{lebo2013prov}
T.~Lebo, S.~Sahoo, D.~McGuinness, K.~Belhajjame, J.~Cheney, D.~Corsar,
D.~Garijo, S.~Soiland-Reyes, S.~Zednik, and J.~Zhao.
\newblock {PROV-O}: The {PROV} ontology.
\newblock {W3C Recommendation}, 2013.

\bibitem{buneman2001and}
P.~Buneman, S.~Khanna, and W.~C.~Tan.
\newblock Why and where: A characterization of data provenance.
\newblock In {\em Proceedings of ICDT}, 2001.

\bibitem{nakamoto2008bitcoin}
S.~Nakamoto.
\newblock Bitcoin: A peer-to-peer electronic cash system.
\newblock 2008.

\end{thebibliography}

\end{document}
